% !TeX program = pdflatex
% !TeX encoding = UTF-8
% Preamble.
\documentclass[12pt]{article}

% Packages.
\usepackage[utf8]{inputenc}     % Unicode support (Umlauts etc.)
\usepackage{amsmath, amssymb}   % Advanced math typesetting
\usepackage{stmaryrd, bbm}
\usepackage[french]{babel}      % Change hyphenation rules
\usepackage[T1]{fontenc}
\usepackage{geometry}
\usepackage{listings}           % Source code formatting and highlighting
\usepackage{csquotes}
\usepackage{interval}
\usepackage{pifont}

% Pages.
\sloppy 
\geometry{a4paper, total={17cm, 25cm}}
\renewcommand{\baselinestretch}{1.4}

%% Real set.
\newcommand{\R}{\mathbb{R}}
%% Length of vector.
\newcommand{\norm}[1]{\left\lVert\ #1\;\right\rVert}
%% Parenthesis.
\newcommand{\of}[1]{\left(#1\right)}

% Data.
\title{Terminale - Mathématiques: géométrie dans l'espace.}
\author{Hector GARLATTI}
\date{\today{}}

\begin{document}
\maketitle

\section{Rappels: vecteurs sur le plan, 2D.}

\paragraph{Notations.}
\begin{itemize}
    \item Le plan en $2$-dimensions est un ensemble: $\R^2$.

    \item Notons un vecteur $\vec{v} \in \R^2$ 
    d'abscisse $x_{\vec{v}} \in \R$ et d'ordonnée: $y_{\vec{v}} \in \R$: $$
        \vec{v} = \begin{pmatrix}
            x_{\vec{v}} \\
            y_{\vec{v}} \\
        \end{pmatrix}
    $$


    \item Notons, aussi, $\overrightarrow{AB}$ le vecteur décrivant le
    \enquote{mouvement} entre le point $A \in \R^2$ et le point $B \in \R^2$: $$
        \overrightarrow{AB} = \begin{pmatrix}
            x_B - x_A \\
            y_B - y_B \\
        \end{pmatrix}
    $$

    \item $
        \norm{\vec{v}}   
        = \sqrt{x_{\vec{v}}^2 + y_{\vec{v}}^2}
        \geq 0
    $: la norme, la longueur du vecteur $\vec{v}$.

    \item Vecteur normal de $\vec{v}$ (notation variable): $$
        \hat{n} = \vec{n} = \hat{v} = \frac{\vec{v}}{\norm{\vec{v}}},
        \quad \norm{\hat{n}} = 1
    $$ 

    \item $A \parallel B$: 
    \begin{itemize}
        \item Si $A$ et $B$ sont droites, segments, courbes:
            $A$ et $B$ sont \textit{parallèles}.
        \item Si $A$ et $B$ sont des vecteurs:
            $A$ et $B$ sont \textit{colinéaires}. 
    \end{itemize}
        
    \item $A \perp B$: 
    \begin{itemize}
        \item Si $A$ et $B$ sont droites, segments, courbes:
            $A$ et $B$ sont \textit{perpendiculaires}.
        \item Si $A$ et $B$ sont des vecteurs:
            $A$ et $B$ sont \textit{orthogonaux}.
    \end{itemize}
    
    \item $H$ est le \textit{projeté orthogonal} du point $A$
        sur $B$ (droite, plan) si: $$
            (AH) \perp B \text{ et } H = (AH) \cap B
        $$
    \item $\widehat{\vec{u}, \vec{v}}$: l'angle orienté entre les vecteurs
        $\vec{u}$ et $\vec{v}$.
\end{itemize}

\paragraph{Géométrie.}
\begin{itemize}
    \item Distance euclidienne: $$
        d = \sqrt{x^2 + y^2}
    $$
    \item Théorème de Thales: soit $2$ triangles $ABC$ et $ADE$: $$
        (BC) \parallel (DE) \iff
        \frac{AB}{AD} = \frac{AC}{AE} = \frac{BC}{DE}
    $$

    \item Milieu $I$ entre deux points $A$ et $B$: $$
        I = \frac{A + B}{2}
    $$

    \item Equation de la droite $(d)$ entre $2$ points $A$ et $B$ de pente $m$. $$
    \begin{aligned}
              &m = \frac{y_B - y_A}{x_B - x_A} \\
        (d):\;&y = m(x - x_A) + y_A = m(x - x_B) + y_B 
    \end{aligned}
    $$
    \item Equation du cercle $\mathcal{C}$ centré sur 
        le point $P$ de rayon $r \in \R$: $$
        \mathcal{C}: (x - x_P)^2 + (y - y_P)^2 = r^2
    $$
\end{itemize}


\paragraph{Vecteurs.}
\begin{itemize}
    \item Négation: $-\overrightarrow{AB} = \overrightarrow{BA}$
    \item Relation de Chasles (\enquote{joindre les bouts}): $
          \overrightarrow{AC} + \overrightarrow{CB}
        = \overrightarrow{AB}
    $
    \item Additions de vecteurs $\vec{u}$ et $\vec{v}$ dans $\R^2$: $$
      \begin{pmatrix} x_u \\ y_u \end{pmatrix}
    + \begin{pmatrix} x_v \\ y_v \end{pmatrix}
    = \begin{pmatrix} x_u + x_v \\ y_u + y_v \end{pmatrix}
    $$
    \item Facteur d'un vecteur $\vec{v} \in \R^2$ par $a \in \R$: $$
          a \times \begin{pmatrix} x_v \\ y_v \end{pmatrix}
        = \begin{pmatrix} a \times x_v \\ a \times y_v \end{pmatrix}
    $$
    \item Produit scalaire de vecteurs $\vec{u}$ et $\vec{v}$ dans $\R^2$: $$
    \begin{aligned}
        \vec{u} \cdot \vec{v} 
        & = x_{\vec{u}} \, x_{\vec{v}} + y_{\vec{u}} \, y_{\vec{v}} \\
        & = \norm{\vec{u}} \times \norm{\vec{v}} \times 
            \cos\of{\widehat{\vec{u}, \vec{v}}}
    \end{aligned}
    $$

\end{itemize}


\section{Géométrie dans un espace euclidien.}
En Terminale, nous étendons notre champ d'action: 
du plan, nous passons à l'espace 3-dimensionnel, dans $\R^3$.
Ainsi, à l'abscisse et à l'ordonné nous ajoutons la \enquote{hauteur}:
la \textit{cote} $z$ et notons, pour $\vec{v} \in \R^3$: $$
    \vec{v} = \begin{pmatrix}
        x_{\vec{v}} \\
        y_{\vec{v}} \\
        z_{\vec{v}} \\
    \end{pmatrix}
$$

Plus formellement, les vecteurs $\vec{i}, \vec{j}, \vec{k}$ 
forment une base $\mathcal{B}$ de $\R^3$ si, et seulement si: $$
    \forall \vec{t} \in \R^3, \quad
    \exists! (a, b, c) \in \R^3 \text{ tel que } 
    \vec{t} = a\vec{i} + b\vec{j} + c\vec{k} 
$$

\paragraph{Propriétés.}
Pour des vecteurs $\vec{u}, \vec{v}, \vec{w}$ dans $\R^3$.

\begin{itemize}
    \item Distance euclidienne dans $\R^3$, sur la même fondation que dans $\R^2$: $$
        d = \sqrt{x^2 + y^2 + z^2}
    $$
    \item $\vec{u} \parallel \vec{v}$: $\vec{u}$ et $\vec{v}$ sont \textit{colinéaires},
        comme dans $\R^2$, quand ils ont la même direction.
        Autrement dit: les droites $(u)$ et $(v)$ guidées par leur vecteur
        respectif sont parallèles $(u) \parallel (v)$. \\
        Formellement, disons: $$
            \exists! \lambda \in \R \quad \text{tel que} \quad
            \vec{u} = \lambda \vec{v}
        $$ \\
        \textsc{Bonus}: nous disons aussi que leur 
        \textit{produit croix} est nul: $\vec{u} \times \vec{v} = 0$

    \item $\vec{u}, \vec{v}, \vec{w}$ sont \textit{coplanaires} s'ils sont compris
        dans une même base $\mathcal{B}$ de $\R^2$. \\
        Formellement, $\vec{u}$ peut s'écrire comme une 
        \textit{combinaison linéaire} de $\vec{v}$ et $\vec{w}$: $$
            \exists(\alpha, \beta) \in (\R^{\star})^2 \quad \text{tel que} \quad
            \vec{u} = \alpha \vec{v} + \beta \vec{w}
        $$
\end{itemize}


\paragraph{Figures.}
\begin{itemize}
    \item Représentation paramétrique de la droite $(d) \subset \R^3$,
        pour un point $P \in (d)$ et un vecteur directeur $\vec{v}$: $$
            (d) \; \begin{cases}
                x = x_P + x_{\vec{v}}\,t \\ 
                y = y_P + y_{\vec{v}}\,t \\ 
                z = z_P + z_{\vec{v}}\,t \\ 
            \end{cases}  \forall t \in \R
        $$
    \item Plan $\mathcal{P} \subset \R^3$: $$
        \mathcal{P}: ax + by + cz + d = 0
    $$
    \item Propriété des milieux, 
        $I \text{ milieu } de \left[AB\right], \forall M \in \R^3$: 
        $$
            2\overrightarrow{MI} = \overrightarrow{MA} + \overrightarrow{MB} 
        $$
    \item $ABCD$ parallélogramme $
        \iff \overrightarrow{AB} + \overrightarrow{AD} = \overrightarrow{AC}
        \iff \overrightarrow{AB} = \overrightarrow{DC}
    $
    \item Sphère $\mathcal{S}$ de rayon $r \in \R_{+}$ et de centre $C \in \R^3$: $$
        \mathcal{S}: (x - x_C)^2 + (y - y_C)^2 + (z - z_C)^2 = r^2
    $$
    
\end{itemize}


\section{Orthogonalité.}

\end{document}
